\capitulo{1}{Introducción}

La enfermedad de Parkinson constituye uno de los desafíos sociosanitarios más urgentes de la actualidad. Clínicamente, se caracteriza por una sintomatología motora progresiva, rigidez y bradicinesia, que compromete severamente la autonomía del individuo. Si bien el tratamiento farmacológico es indispensable, existe un consenso médico absoluto sobre el papel de la rehabilitación física continuada como el único mecanismo eficaz para preservar la independencia funcional del paciente a largo plazo.

Esta necesidad es aún más urgente si recordamos que la enfermedad de Parkinson afecta mayoritariamente a personas de edad avanzada. Para este colectivo, la rehabilitación no es solo un ejercicio físico, sino la única herramienta que les permite seguir valiéndose por sí mismos en su día a día.

Sin embargo, la eficacia clínica de la rehabilitación choca frontalmente con una barrera logística y demográfica: la inequidad en el acceso al tratamiento. En un contexto marcado por el envejecimiento poblacional y la dispersión geográfica —fenómenos especialmente agudos en los territorios de la denominada ``España Vaciada''—, el modelo tradicional de rehabilitación presencial se vuelve insostenible.

Para un paciente con movilidad reducida que reside en un entorno rural, el desplazamiento diario a un centro hospitalario de referencia no supone únicamente un esfuerzo físico, sino una carga logística y económica que, a menudo, recae sobre las familias. Esta fricción provoca que, en la práctica, muchos pacientes no reciban la terapia necesaria. Ante la saturación de los servicios sanitarios y la dificultad de cobertura en zonas aisladas, surge la necesidad imperiosa de desarrollar soluciones tecnológicas que trasladen la terapia del hospital al domicilio, garantizando la continuidad asistencial sin perder la supervisión facultativa.

En este escenario, los dispositivos móviles se presentan como la herramienta ideal para acercar la terapia al paciente. Su uso se ha extendido enormemente en los últimos años, convirtiéndose en un elemento cotidiano incluso para las personas mayores. Además de su popularidad, ofrecen una ventaja fundamental: su interacción es mucho más intuitiva. Gracias a las pantallas táctiles, se eliminan las barreras que suelen tener los ordenadores convencionales, permitiendo que el paciente acceda a sus ejercicios de forma natural y sencilla.

\section{Contexto académico}

El presente Trabajo de Fin de Grado no se constituye como una iniciativa aislada, sino que se enmarca en una línea de investigación activa de la Universidad de Burgos. El propósito superior de esta investigación es evaluar el impacto de la tecnología en la mejora de la calidad de vida de pacientes con enfermedades neurodegenerativas.

Dentro de este marco, el proyecto se centra en el desarrollo de una infraestructura de telerrehabilitación basada en el uso de dispositivos móviles. El objetivo principal ha sido construir el sistema necesario para capturar y digitalizar los movimientos del paciente a través de la cámara del teléfono o la tableta. Esta herramienta sienta las bases para crear modelos de comparación y evaluación de ejercicios que permitan una telerrehabilitación totalmente autónoma sin la necesidad de un terapeuta.

\section{Justificación y motivación del proyecto}

La elección de este proyecto responde a un doble objetivo: por un lado, aplicar la ingeniería para resolver un problema social real y, por otro, afrontar el reto técnico que supone desarrollar una aplicación completa.

La motivación primaria reside en la vocación de utilidad pública. Frente a la opción de realizar un ejercicio académico teórico, se optó por abordar una problemática tangible: la brecha de acceso a la salud en entornos rurales. El proyecto se concibe como una herramienta de \textit{e-Health} destinada a democratizar la rehabilitación, permitiendo que un paciente de un núcleo rural acceda a la misma calidad de seguimiento que uno urbano. La posibilidad de aportar una solución técnica que impacte directamente en la calidad de vida de un colectivo vulnerable dota al trabajo de una dimensión ética que trasciende la mera programación.

Desde la perspectiva profesional, este TFG se planteó como el escenario idóneo para la transición hacia la ingeniería de software profesional. Al requerir el desarrollo de una solución integral, abarcando desde el diseño de bases de datos relacionales y APIs en el servidor hasta la gestión de hardware multimedia y concurrencia en el dispositivo móvil, el proyecto exigió abandonar la zona de confort académica. Esta complejidad obligó a adoptar un enfoque de autoaprendizaje para dominar tecnologías de desarrollo nativo y arquitecturas modernas demandadas en el mercado laboral actual.

\section{Objetivos generales}

La finalidad última del sistema es eliminar la barrera de la distancia física mediante un puente digital. Se ha diseñado y desarrollado una aplicación móvil nativa que satisface dos funciones críticas:

\begin{itemize}
    \item \textbf{Accesibilidad universal para el paciente:} proveer una herramienta de ejecución de ejercicios guiada por vídeo. Dada la afectación motora y la brecha digital del usuario final, el objetivo prioritario ha sido diseñar una interfaz adaptada a los temblores y la edad, y extremadamente sencilla, garantizando la adherencia al tratamiento.
    
    \item \textbf{Control asíncrono para el terapeuta:} dotar al profesional sanitario de una plataforma de gestión eficiente. El sistema permite la asignación remota de terapias y, fundamentalmente, la supervisión diferida (asíncrona) mediante la revisión de los vídeos grabados por el paciente, optimizando los tiempos de consulta y permitiendo escalar la atención a un mayor número de usuarios.
\end{itemize}

\section{Entregables y repositorio}

El resultado tangible de este trabajo es un prototipo funcional de aplicación Android que integra los roles de Administrador, Terapeuta y Paciente, cubriendo el ciclo completo de rehabilitación: prescripción, ejecución, grabación y evaluación.

Para facilitar la auditoría técnica y la reproducibilidad del proyecto, la totalidad del código fuente, junto con la documentación técnica y los scripts de despliegue, se encuentra disponible en el siguiente repositorio de control de versiones:

\begin{center}
    \url{https://github.com/andrespuentesg/TFG-GII24.36-AndresPuentes}
\end{center}